\acknowledgments
\label{sec:acknowledgements}

\hspace{0.85cm}JCdJ acknowledges support from Fondation CentraleSupélec and Fondation Ailes de France. BD acknowledges support from the Ambrose Monell Foundation, the Corning Glass Works Foundation Fellowship Fund, and the Institute for Advanced Study. SGG acknowledges that this work was performed in part at the Aspen Center for Physics, which is supported by National Science Foundation grant PHY-2210452 and a grant from the Alfred P Sloan Foundation (G-2024-22395). %SGG acknowledges funding from the Virginia Institute for Theoretical Astrophysics (VITA), supported by the College and Graduate School of Arts and Sciences at the University of Virginia. 
US is supported by NSF CDSE grant number AST-2408026 and NASA TCAN grant number 80NSSC24K0101. TZ is supported by Schmidt Sciences.
 
%This research used data obtained with the Dark Energy Spectroscopic Instrument (DESI). DESI construction and operations is managed by the Lawrence Berkeley National Laboratory. This material is based upon work supported by the U.S. Department of Energy, Office of Science, Office of High-Energy Physics, under Contract No. DE–AC02–05CH11231, and by the National Energy Research Scientific Computing Center, a DOE Office of Science User Facility under the same contract. Additional support for DESI was provided by the U.S. National Science Foundation (NSF), Division of Astronomical Sciences under Contract No. AST-0950945 to the NSF’s National Optical-Infrared Astronomy Research Laboratory; the Science and Technology Facilities Council of the United Kingdom; the Gordon and Betty Moore Foundation; the Heising-Simons Foundation; the French Alternative Energies and Atomic Energy Commission (CEA); the National Council of Humanities, Science and Technology of Mexico (CONAHCYT); the Ministry of Science and Innovation of Spain (MICINN), and by the DESI Member Institutions: \url{www.desi.lbl.gov/collaborating-institutions}. The DESI collaboration is honored to be permitted to conduct scientific research on I’oligam Du’ag (Kitt Peak), a mountain with particular significance to the Tohono O’odham Nation. Any opinions, findings, and conclusions or recommendations expressed in this material are those of the author(s) and do not necessarily reflect the views of the U.S. National Science Foundation, the U.S. Department of Energy, or any of the listed funding agencies.

The Hyper Suprime-Cam (HSC) collaboration includes the astronomical communities of Japan and Taiwan, and Princeton University. The HSC instrumentation and software were developed by the National Astronomical Observatory of Japan (NAOJ), the Kavli Institute for the Physics and Mathematics of the Universe (Kavli IPMU), the University of Tokyo, the High Energy Accelerator Research Organization (KEK), the Academia Sinica Institute for Astronomy and Astrophysics in Taiwan (ASIAA), and Princeton University. Funding was contributed by the FIRST program from the Japanese Cabinet Office, the Ministry of Education, Culture, Sports, Science and Technology (MEXT), the Japan Society for the Promotion of Science (JSPS), Japan Science and Technology Agency (JST), the Toray Science Foundation, NAOJ, Kavli IPMU, KEK, ASIAA, and Princeton University. This paper makes use of software developed for Vera C. Rubin Observatory. We thank the Rubin Observatory for making their code available as free software at \url{http://pipelines.lsst.io/}. This paper is based on data collected at the Subaru Telescope and retrieved from the HSC data archive system, which is operated by the Subaru Telescope and Astronomy Data Center (ADC) at NAOJ. Data analysis was in part carried out with the cooperation of Center for Computational Astrophysics (CfCA), NAOJ. We are honored and grateful for the opportunity of observing the Universe from Maunakea, which has the cultural, historical and natural significance in Hawaii. 

This analysis made use of the following software packages: \texttt{CosmoSIS} (\url{https://cosmosis.readthedocs.io/en/latest/}) \cite{CosmoSISZuntz2015}, \texttt{pocoMC} (\url{https://pocomc.readthedocs.io/en/latest/}) \cite{karamanis2022pocoMC1, karamanis2022pocoMC2}, \texttt{getDist}(\url{https://getdist.readthedocs.io/en/latest/})\cite{getDistLewis2025}, \texttt{ChainConsumer} (\url{https://github.com/Samreay/ChainConsumer}) \cite{ChainConsumer_Hinton2016}, \texttt{matplotlib} \cite{Hunter2007Matplotlib}, \texttt{numpy} \cite{harris2020numpy}.
%This analysis used the following software packages: \texttt{astropy} \cite{astropy2013paper1, astropy2018paper2, astropy2023paper3},  \texttt{CCL} \cite{CCLChisari2019}, \texttt{Corrfunc}  \cite{Sinha2019corrfunc1, Sinha2019corrfunc2} (with \texttt{cosmodesi/pycorr}), \texttt{mocpy} \cite{fernique2015moc},  \texttt{numpy} \cite{harris2020numpy}, \texttt{SciPy} \cite{2020SciPyNMeth}, \texttt{matplotlib} \cite{Hunter2007Matplotlib}, \texttt{pyMC} \cite{Abril-Pla_PyMC_a_modern_2023}, \texttt{scikit-learn}  \cite{ScikitLearnPedregosa2011}, \texttt{fitsio}.